% ============================================================
% 2. Related Work
% ============================================================

\section{Related Work}
\label{sec:related}

\paragraph{Multi-agent LLM frameworks.}
LangGraph~\citep{langgraph2024} models agent workflows as state
machines with checkpointing but provides no memory bound.
AutoGen~\citep{wu2023autogen} enables flexible multi-agent
conversations; CrewAI~\citep{crewai2024} and
CAMEL~\citep{li2023camel} orchestrate role-based agent teams;
MetaGPT~\citep{hong2023metagpt} assigns software-engineering roles.
All store every agent state in unbounded Python data structures.
Our work is orthogonal: BIC can serve as the memory layer beneath any
of these frameworks.

\paragraph{Memory management for LLM agents.}
MemGPT~\citep{packer2023memgpt} (now Letta) introduces an
OS-inspired paging system for a \emph{single} agent's context window,
moving segments to disk when the context is full.  It does not address
multi-agent hierarchies or provide formal memory bounds.  We
implement a Tiered baseline mirroring MemGPT's hot/cold architecture
(\S\ref{sec:stress}); BIC outperforms it by $32\times$ on
reconstruction quality at scale.
Reflexion~\citep{shinn2023reflexion} compresses a single agent's
episodic memory into a verbal summary for the next trial.
Generative Agents~\citep{park2023generative} maintains an
ever-growing memory stream with retrieval over it.
None of these systems bound total memory across an agent swarm or
provide reconstruction guarantees for evicted states.

\paragraph{Recent memory management for LLM agents.}
The growing interest in long-running LLM agents has spurred recent work
on agent memory management.  The ICLR 2026 MemAgents
Workshop~\citep{memagents2026} brought together researchers addressing
memory architectures, retrieval mechanisms, and resource management for
LLM-based agents, highlighting the timeliness of this problem.
A-MEM~\citep{amem2025} proposes an agentic memory system where an LLM
agent autonomously decides what to remember and how to organize its
memories, treating memory management itself as an agentic task.
Both lines of work focus on the \emph{quality} and \emph{organization}
of individual agent memories, whereas BIC addresses a complementary
\emph{systems} challenge: bounding the total cache memory consumed by
an entire swarm of agents while preserving retrievability of evicted
states through hierarchical summarization.  BIC's cache management
layer could be combined with richer per-agent memory strategies such
as those explored in A-MEM.

\paragraph{Streaming and sketching algorithms.}
The streaming-algorithms community provides tools for processing
unbounded data in bounded space: Count-Min
Sketches~\citep{cormode2005cm}, HyperLogLog~\citep{flajolet2007hyperloglog},
and sliding-window algorithms~\citep{datar2002sliding}.  These
structures are designed for aggregate statistics (counts, cardinalities)
over data streams, not for managing structured agent states with
parent-child relationships and reconstruction requirements.  BIC
borrows the principle of bounded-space processing but applies it to
hierarchical agent-state management with summarization.

\paragraph{Cache-efficient computation and space-filling curves.}
Hilbert curves are widely used in database
indexing~\citep{lawder2000hilbert} and cache-oblivious
algorithms~\citep{frigo1999cache} to improve spatial
locality.  The Cantor pairing function provides a bijection
$\Nat \times \Nat \to \Nat$ used in computability theory and database
design.  We repurpose both for a novel application: mapping
tree-structured agent hierarchies to cache slots (Cantor) and
clustering agent states for locality-aware eviction (Hilbert).

\paragraph{Hierarchical summarization.}
Tree-based summarization appears in document
summarization~\citep{wu2021recursively} and retrieval-augmented
generation~\citep{raptor2024}.  Our hierarchical summarization differs
in that it is \emph{eviction-driven}: summaries are created lazily
when cache pressure forces eviction, and the summarization chain
doubles as the reconstruction path for queries about evicted agents.
